\section{Rows to columns}\label{sec:demand}

Theorems~H, X and~S use transposition only through the certificate method: a transposer cannot have a large acyclic set after
a cheap cut. This section uses the transposition demand directly. Information from a set of input rows must reach the output rows
into which transposition sends it, and every cut between the two must carry it.

For sets $A_0,A_1\subseteq\{0,\dots,d-1\}$ of rows let $\Lin(A_0)$ be the set of steps at which the input head reads a position in
an input row of $A_0$ (\lean{inSlabReadSteps}), and let $\Qstar(A_0,A_1)$ be the set of input positions whose input row lies in
$A_0$ and whose output row lies in $A_1$ (\lean{bichromaticRect}): the bits that rows $A_0$ send into output rows $A_1$. It has
exactly $\abs{A_0}\abs{A_1}b$ elements (\lean{card_bichromaticRect}).

\begin{lemma}[Lemma~R, rows to columns]\label{lem:R}
\formal{TranspositionCut.lean}{isAcyclicIn_coneSystem_uncutRootedBichromaticRect, frozenTrace_bichromatic_tposOf_cut_ge, oblivious_bichromatic_tposOf_cut_ge}
Let $M$ transpose at size $d\ge1$ within time $c$, and let $\tau$ be an $\eta$-heavy trace of $M$; for an oblivious $M$ take its common
trace and $\eta=0$. For all sets $A_0,A_1$ of rows and every set $\Gamma$ of steps separating $\Out(\Qstar(A_0,A_1))$ from $\Lin(A_0)$,
\[
  \wval\abs\Gamma+\abs{\Gamma\cap\Out(\Qstar(A_0,A_1))}\ \ge\ \abs{A_0}\abs{A_1}b-\eta L,
  \qquad\text{hence}\qquad \abs\Gamma\ \ge\ \frac{\abs{A_0}\abs{A_1}b-\eta L}{\wval+1}.
\]
\end{lemma}

\begin{outline}
Let $Q'$ be the set of $q\in\Qstar(A_0,A_1)$ whose root is not in $\Gamma$; since $\Out$ is injective,
$\abs{Q'}\ge\abs{A_0}\abs{A_1}b-\abs{\Gamma\cap\Out(\Qstar)}$. For $q\in Q'$, a dependency path from $\Out(q)$ that avoids $\Gamma$ after
its root cannot end at a read of rows $A_0$: the root is not in $\Gamma$, so the path would avoid $\Gamma$ altogether. Hence the cone of
every $q\in Q'$ avoiding $\Gamma$ contains no position of rows $A_0$, while $Q'$ lies in rows $A_0$. No element of $Q'$ sees another, so
$Q'$ is acyclic, and \cref{lem:frozen} gives $\abs{Q'}-\wval\abs\Gamma\le\eta L$.
\end{outline}

With $A_0=A_1$ the set of all rows, Lemma~R says that every cut between the input and the output carries the whole input:
$\abs\Gamma\ge(1-\eta)L/(\wval+1)$. With smaller row sets it says the same for every block of the transposition demand separately. The
filtering of roots matters: $\Gamma=\Out(\Qstar)$ is a separating set, and the cone relation avoiding it need not be acyclic on all of
$\Qstar$, because roots in $\Gamma$ may still see their own inputs.

The lemma becomes a statement about a single set of steps through a colouring.

\begin{theorem}[Colouring criterion]\label{thm:colouring}
\formal{TranspositionCut.lean}{isMengerVertexCut_of_bichromaticColoring, frozenTrace_bichromaticColoring_wave_cut_ge}
Let $M$, $d$, $c$, $\eta$ and $\tau$ be as in \cref{lem:R}, with halting time $T$, and let $A_0,A_1$ be sets of rows. Let $Y$ be a set of
steps and $\chi$ a colouring of the steps by $\{0,1\}$ such that
\begin{enumerate}[label=(\roman*)]
\item $\chi=0$ on $\Lin(A_0)\setminus Y$ and $\chi=1$ on $\Out(\Qstar(A_0,A_1))\setminus Y$;
\item every revisit $(u,t)$ with $u,t\notin Y$ has $\chi(t)\le\chi(u)$.
\end{enumerate}
Let $\partial\chi$ be the set of steps $t\in[1,T)$ with $t,t-1\notin Y$, $\chi(t)=1$ and $\chi(t-1)=0$. Then every set $B$ of steps separating
the roots of all input positions from $Y$ satisfies
\[
  \abs B\ \ge\ \frac{\abs{A_0}\abs{A_1}b-\eta L}{\wval+1}-\abs{\partial\chi}.
\]
\end{theorem}

\begin{outline}
$B\cup\partial\chi$ separates $\Out(\Qstar)$ from $\Lin(A_0)$. Take a dependency path from a root to a read of rows $A_0$ that avoids
$B\cup\partial\chi$. It avoids $Y$, since its part up to a vertex of $Y$ would be a path from a root to $Y$ avoiding $B$. By (i) the path
starts with colour $1$ and ends with colour $0$, so some edge drops from colour $1$ to colour $0$ as the path moves to an earlier step. A
state edge $t\to t-1$ with $\chi(t)=1$ and $\chi(t-1)=0$ would put $t$ in $\partial\chi$, and a revisit edge $t\to u$ with $\chi(t)=1$ and
$\chi(u)=0$ is excluded by (ii), a contradiction. \Cref{lem:R} bounds $\abs{B\cup\partial\chi}$.
\end{outline}

Taking $B=Y$ shows $\abs Y\ge(\abs{A_0}\abs{A_1}b-\eta L)/(\wval+1)-\abs{\partial\chi}$: a set of steps across which the colour can only
jump at few places must itself be large. Several such sets that are pairwise disjoint give a time bound. For $d\ge16$ let $K=\lfloor
\log_2d\rfloor$ and $d_1=2^K$, and for $0\le s\le K-2$ let $A_s^0$ and $A_s^1$ be the rows $i<d_1$ whose bit $K-1-s$ is $0$, respectively $1$.

\begin{corollary}[Scale-separated traces]\label{cor:scale-separated}
\notformalized
Let $M$, $d$, $c$, $\eta$ and $\tau$ be as in \cref{lem:R}, with $d\ge16$, $\eta\le\frac1{32}$, $K_{\mathrm{flip}}\ge0$ and
$d\ge64(\wval+1)K_{\mathrm{flip}}$. Suppose that for each $s$ in a set $S\subseteq\{0,\dots,K-2\}$ there are a set $Y_s$ of steps and a
colouring $\chi_s$ satisfying (i) and (ii) of \cref{thm:colouring} for $(A_s^0,A_s^1)$ with $\abs{\partial\chi_s}\le K_{\mathrm{flip}}d$, and
that the sets $Y_s$ are pairwise disjoint. Then $T\ge\abs S\cdot L/(64(\wval+1))$.
\end{corollary}

\begin{informalproof}
Since $d_1>d/2$, $\abs{A_s^0}\abs{A_s^1}b=(d_1/2)^2b>L/16$. Since $d\ge64(\wval+1)K_{\mathrm{flip}}$ and $b\ge1$,
$K_{\mathrm{flip}}d\le d^2/(64(\wval+1))\le L/(64(\wval+1))$. \Cref{thm:colouring} with $B=Y_s$ gives
\[
  \abs{Y_s}\ \ge\ \frac{(1/16-\eta)L}{\wval+1}-K_{\mathrm{flip}}d\ \ge\ \frac{L}{32(\wval+1)}-\frac{L}{64(\wval+1)}=\frac{L}{64(\wval+1)} .
\]
The $Y_s$ are pairwise disjoint subsets of $\Range T$, so $T\ge\sum_{s\in S}\abs{Y_s}$.
\end{informalproof}

With $S$ all $K-1$ scales the bound is $T\ge(\lfloor\log_2d\rfloor-1)L/(64(\wval+1))$, so no linear-time transposer has a trace that is
scale-separated across all scales. The corollary should be read with \cref{prop:onecut} in mind: one set of steps, together with the
colouring that changes once, at the step after the last input read, satisfies (i) and (ii) for every pair of row sets at once, so the content
of the hypothesis is the disjointness of the sets $Y_s$, which says roughly that live storage is rewritten once per scale. The corollary is
therefore a time bound for schedules that refine the row index one bit at a time in separate passes, whatever they compute; it does not
show that a linear-time transposer must do so.
